\documentclass[11pt]{article}
\usepackage[a4paper,margin=0.8in,top=0.85in,bottom=0.9in]{geometry}
\usepackage[T1]{fontenc}
\usepackage{charter}
\usepackage[scaled=0.92]{helvet}
\usepackage{amsmath,amssymb}
\usepackage{booktabs}
\usepackage{tikz}
\usepackage{titlesec}
\usepackage{xcolor}
\usepackage{enumitem}
\usepackage{caption}
\usepackage[hidelinks]{hyperref}

% \num is used only to keep plain numbers out of the prose; no
% package needed for something this simple.
\newcommand{\num}[1]{#1}

% Generated by src/study.py - do not edit by hand.
% Rerun 'python3 run.py' to refresh these numbers.
\newcommand{\NLocations}{32}
\newcommand{\NComplaintSites}{31}
\newcommand{\NDorms}{21}
\newcommand{\NNewCampus}{24}
\newcommand{\NOldCampus}{8}
\newcommand{\NTunnelMetres}{94.3}
\newcommand{\NMeanWalk}{5.7}
\newcommand{\NMeanSameCampus}{2.1}
\newcommand{\NMeanCrossCampus}{11.3}
\newcommand{\NMaxWalk}{16.1}
\newcommand{\NFifoObj}{3602.6}
\newcommand{\NFifoTard}{3241.0}
\newcommand{\NFifoTravel}{26.8}
\newcommand{\NFifoSla}{85.1}
\newcommand{\NFifoUrgent}{37.0}
\newcommand{\NFifoUrgentWait}{4.47}
\newcommand{\NFifoUrgentNinety}{7.26}
\newcommand{\NFifoBacklog}{6.1}
\newcommand{\NFifoClosed}{1657}
\newcommand{\NFifoUrgentLate}{153}
\newcommand{\NFifoUrgentUnavoidable}{32}
\newcommand{\NFifoUrgentAvoidable}{121}
\newcommand{\NSavObj}{2766.2}
\newcommand{\NSavTard}{2563.8}
\newcommand{\NSavTravel}{15.0}
\newcommand{\NSavSla}{90.5}
\newcommand{\NSavUrgent}{55.5}
\newcommand{\NSavUrgentWait}{3.47}
\newcommand{\NSavUrgentNinety}{6.33}
\newcommand{\NSavBacklog}{3.3}
\newcommand{\NSavClosed}{1665}
\newcommand{\NSavUrgentLate}{109}
\newcommand{\NSavUrgentUnavoidable}{32}
\newcommand{\NSavUrgentAvoidable}{77}
\newcommand{\NPriObj}{358.0}
\newcommand{\NPriTard}{106.5}
\newcommand{\NPriTravel}{18.6}
\newcommand{\NPriSla}{97.9}
\newcommand{\NPriUrgent}{85.7}
\newcommand{\NPriUrgentWait}{2.92}
\newcommand{\NPriUrgentNinety}{4.07}
\newcommand{\NPriBacklog}{0.2}
\newcommand{\NPriClosed}{1665}
\newcommand{\NPriUrgentLate}{35}
\newcommand{\NPriUrgentUnavoidable}{32}
\newcommand{\NPriUrgentAvoidable}{3}
\newcommand{\NImpObj}{319.9}
\newcommand{\NImpTard}{108.3}
\newcommand{\NImpTravel}{15.7}
\newcommand{\NImpSla}{98.0}
\newcommand{\NImpUrgent}{86.9}
\newcommand{\NImpUrgentWait}{2.42}
\newcommand{\NImpUrgentNinety}{4.05}
\newcommand{\NImpBacklog}{0.6}
\newcommand{\NImpClosed}{1665}
\newcommand{\NImpUrgentLate}{32}
\newcommand{\NImpUrgentUnavoidable}{32}
\newcommand{\NImpUrgentAvoidable}{0}
\newcommand{\NUrgentTotal}{245}
\newcommand{\NUrgentImpossible}{32}
\newcommand{\NTravelCutPct}{41.4}
\newcommand{\NObjCutPct}{91.1}
\newcommand{\NWalkSavedPerDay}{100.0}
\newcommand{\NBbMaxTickets}{10}
\newcommand{\NBbGapWorst}{0.0}
\newcommand{\NBbServed}{7}
\newcommand{\NBbNodesArbCold}{22,721}
\newcommand{\NBbNodesArbWarm}{10,131}
\newcommand{\NBbOrderCutPct}{52.9}
\newcommand{\NBbWarmCutPct}{5.3}
\newcommand{\NBbNodesWarm}{10,131}
\newcommand{\NBbNodesCold}{10,700}
\newcommand{\NBbNodeCutPct}{5.0}
\newcommand{\NBbSeconds}{0.02}
\newcommand{\NBbSavWorstGap}{4342.5}
\newcommand{\NBbSavWorstJobs}{8}
\newcommand{\NBbSavWorstTimes}{44.4}
\newcommand{\NCrewSixObj}{2717.6}
\newcommand{\NCrewSixSla}{90.3}
\newcommand{\NCrewSixUrgent}{79.6}
\newcommand{\NCrewNineObj}{319.9}
\newcommand{\NCrewNineSla}{98.0}
\newcommand{\NCrewNineUrgent}{86.9}
\newcommand{\NCrewTwelveObj}{293.1}
\newcommand{\NCrewTwelveSla}{98.1}
\newcommand{\NCrewTwelveUrgent}{86.9}
\newcommand{\NCrewRebalObj}{2355.7}
\newcommand{\NCrewRebalSla}{93.4}
\newcommand{\NCrewRebalUrgent}{85.3}
\newcommand{\NSurgeLowFifoSla}{85.1}
\newcommand{\NSurgeLowFifoBacklog}{6.1}
\newcommand{\NSurgeLowFifoUrgent}{37.0}
\newcommand{\NSurgeLowImpSla}{98.0}
\newcommand{\NSurgeLowImpBacklog}{0.6}
\newcommand{\NSurgeLowImpUrgent}{86.9}
\newcommand{\NSurgeHighFifoSla}{53.9}
\newcommand{\NSurgeHighFifoBacklog}{113.5}
\newcommand{\NSurgeHighFifoUrgent}{18.6}
\newcommand{\NSurgeHighImpSla}{88.1}
\newcommand{\NSurgeHighImpBacklog}{23.8}
\newcommand{\NSurgeHighImpUrgent}{83.9}
\newcommand{\NSeedCount}{5}
\newcommand{\NSeedImpMin}{277.9}
\newcommand{\NSeedImpMax}{348.5}
\newcommand{\NSeedFifoMin}{3562.4}
\newcommand{\NSeedFifoMax}{4825.5}
\newcommand{\NDayTickets}{58}
\newcommand{\NDayFifoTravel}{239.0}
\newcommand{\NDayImpTravel}{160.9}
\newcommand{\NDays}{30}
\newcommand{\NTicketsPerDay}{58}
\newcommand{\NWorkingMinutes}{450}
\newcommand{\NAlpha}{0.05}
\newcommand{\NAlphaZeroTravel}{18.3}
\newcommand{\NAlphaWorthMinutes}{2.6}

\definecolor{figink}{HTML}{1A1A1A}
\definecolor{figmuted}{HTML}{6B7280}
\definecolor{figfaint}{HTML}{D1D5DB}
\definecolor{figpaper}{HTML}{FFFFFF}
\definecolor{figaccent}{HTML}{1F4E79}
\definecolor{figdorm}{HTML}{2E7D32}
\definecolor{figoffice}{HTML}{6A4C93}
\definecolor{figfacility}{HTML}{B26500}
\definecolor{figdepot}{HTML}{C62828}
\definecolor{figroad}{HTML}{9AA3AC}
\definecolor{figtunnel}{HTML}{1F4E79}
\definecolor{figr1}{HTML}{1F4E79}
\definecolor{figr2}{HTML}{C62828}
\definecolor{figr3}{HTML}{2E7D32}
\definecolor{figr4}{HTML}{B26500}
\definecolor{figr5}{HTML}{6A4C93}
\definecolor{figr6}{HTML}{00796B}
\definecolor{figr7}{HTML}{AD1457}
\definecolor{figr8}{HTML}{5D4037}
\definecolor{figr9}{HTML}{37474F}

\definecolor{ink}{HTML}{1A1A1A}
\definecolor{accent}{HTML}{1F4E79}
\definecolor{rulegrey}{HTML}{B0B8C0}
\color{ink}

\titleformat{\section}{\normalfont\bfseries\sffamily\color{accent}\large}{\thesection.}{0.6em}{}
\titlespacing{\section}{0pt}{16pt}{5pt}
\titleformat{\subsection}{\normalfont\bfseries\sffamily\color{ink}\normalsize}{\thesubsection}{0.6em}{}
\titlespacing{\subsection}{0pt}{11pt}{3pt}

\captionsetup{font=small,labelfont={bf,sf},labelsep=period,width=0.92\textwidth}
\setlist[itemize]{leftmargin=1.3em,itemsep=2pt,topsep=3pt}
\setlength{\parskip}{4pt}
\setlength{\parindent}{0pt}
\linespread{1.04}

\newcommand{\figbox}[2]{%
  \begin{figure}[t]\centering
  \resizebox{\linewidth}{!}{\input{../out/figures/#1.tikz}}
  \caption{#2}\end{figure}}

\begin{document}

\begin{center}
{\sffamily\bfseries\LARGE\color{accent} Dispatching the dorm maintenance crew}\\[5pt]
{\sffamily\large An optimisation model for the IIMA Estate Office}\\[7pt]
{\small Operations Management \& Logistics \quad\textbar\quad group project \quad\textbar\quad
built on the real campus map, run on synthetic complaint data}
\end{center}
\vspace{2pt}
{\color{rulegrey}\hrule height 0.7pt}
\vspace{6pt}

{\small\itshape Every number in this report is produced by the code that
accompanies it. Running \texttt{python3 run.py} regenerates the results, the
figures and the numbers quoted in the text, so nothing here can drift away
from the model it describes.}

\section{The problem}

The Estate Office receives maintenance complaints from all over campus, all
day long. A tubelight has died in one dorm, a tap is leaking in another, a
window will not shut in a third, somebody's switchboard is sparking. Each
complaint has to be given to a technician, and that technician has to walk to
the room, do the repair, and move on to the next job. A single technician
might get through six or eight jobs in a day.

At the moment the queue is worked more or less in the order the complaints
arrive. That sounds fair, and in a way it is, but it produces two problems
that anybody who has waited a week for a fan to be fixed will recognise.

The first is wasted walking. Complaints arrive in no particular geographic
order, so working the queue in sequence sends a technician from one end of
campus to the other and back again. This matters more at IIMA than it would
elsewhere, because the campus is cut in two by a main road. The Estate Office
sits on the old campus, while \NDorms{} of the \NComplaintSites{} places
that raise complaints are dormitories on the new campus. There is no crossing at
road level. The only way across on foot is the walkway tunnel, which means
every trip between the two halves costs about \NMeanCrossCampus{} minutes,
against roughly \NMeanSameCampus{} minutes for a trip that stays on one side.
A plan that crosses the road four times has thrown away the better part of an
hour before anybody has picked up a screwdriver.

The second problem is that first-come-first-served ignores urgency entirely.
A blocked drain and a squeaky hinge are treated identically if they were
reported in that order. Over the \NDays{} simulated days in this study, the
first-come-first-served baseline closes only \NFifoUrgent\% of urgent
complaints within the promised four hours.

This project asks a simple question: given the complaints waiting at the
start of a day, which technician should take which jobs, and in what order?

\section{Why this is harder than it looks}

Four things make the problem difficult, and they interact.

\textbf{It is a routing problem, so the number of possible answers explodes.}
A single technician with eight jobs can do them in $8! = 40{,}320$ different
orders. Deciding which of nine technicians takes which of sixty jobs, and then
sequencing each of them, is a far larger problem again. This is the vehicle
routing problem, with technicians in place of vehicles and the Estate Office
as the depot, and it is NP-hard. There is no practical way to check every
possibility, so beyond a handful of jobs we have to use heuristics and then
work out how good they are.

\textbf{The thing we care about is not distance.} A delivery company
minimises kilometres. We care about complaints being closed on time, and
walking only matters because time spent walking is time not spent repairing.
These two goals genuinely conflict. Bundling all the jobs in one dorm into a
single visit is excellent for walking and can be terrible for an urgent
complaint in a dorm nobody is visiting until four o'clock.

\textbf{Access windows break a convenient property.} A dorm room needs its
occupant there to open the door, and students are in class for much of the
day. Once jobs have windows, the cost of a route depends on the \emph{order}
of the jobs and not merely on which jobs are in it. That is what makes the
problem awkward to solve: a merge or a reordering that looks sensible on the
map can make a route infeasible, so every candidate has to be re-checked
rather than assumed valid.

\textbf{The geography is not what it looks like.} Two dorms a hundred metres
apart on opposite sides of the road are a twelve minute walk apart. Any model
that uses straight-line distances will quietly produce plans that cannot be
walked.

\section{The campus, and how we measured walking}

We took the campus layout from OpenStreetMap rather than inventing it. The
extract gives real coordinates for \NLocations{} places that matter here:
\NDorms{} dormitories, the academic and administrative blocks, the library,
the dining hall and the student activity block, and the Estate Office itself,
which we use as the depot. \NNewCampus{} of them are on the new campus and
\NOldCampus{} on the old campus. The walkway tunnel is in the map data too,
all \NTunnelMetres{} metres of it.

Walking times are then built in a way that keeps the one feature that really
matters. We treat the two sides of the road as separate, and we join them only
through the tunnel. Within one side, the walking distance between two places
is their straight-line distance multiplied by a detour factor of
\num{1.3}, because paths go around lawns and buildings rather than through
them. Each place is also connected to the tunnel mouth on its own side.
Dijkstra's algorithm run over this small network then works out, for every
pair of places, whether the trip stays on one side or has to go under the
road, and what it costs.

The result behaves the way the campus does. The average walk between two
places on the same side is \NMeanSameCampus{} minutes. The average walk
between two places on opposite sides is \NMeanCrossCampus{} minutes. The
longest trip on campus is \NMaxWalk{} minutes. Because the depot is on the
old campus and almost all the complaints are on the new campus, every
technician starts the day on the wrong side of the road, which is exactly the
situation that makes bundling worthwhile.

\figbox{campus}{The campus as the model sees it. The Estate Office, in red,
is the depot: every technician starts and ends the day there. The nineteen
dormitories are the green cluster on the new campus. Nothing crosses the main
road except the walkway tunnel, so every trip between the two halves is
funnelled through it.}

\section{The complaints}

Nobody was going to hand us the Estate Office's real ticket log, so we
generate complaints ourselves. The point of doing it this way is not to guess
the true numbers. It is that every assumption is written down in one file and
can be changed in one place, so anybody who does have the real figures can put
them in and rerun the study. Table~\ref{tab:assume} lists what we assumed.

\begin{table}[h]\centering\small
\begin{tabular}{llp{58mm}}
\toprule
\textbf{What} & \textbf{Value we used} & \textbf{Where it comes from} \\
\midrule
Working day & 450 minutes & 09:00 to 17:30 with an hour for lunch \\
Walking speed & 75 m/min & 4.5 km/h, an unhurried pace with a toolbag \\
Detour factor & 1.3 & paths go around lawns, not through them \\
Tunnel penalty & 2.0 minutes & down the ramp and back up \\
Complaints per day & 58 on average & assumed; varied in Section 9 \\
Share from dorms & 85\% & dorms hold most of the people \\
Trade mix & 45 / 35 / 20\% & electrical, plumbing, carpentry \\
Repair time & 22 / 38 / 30 min typical & by trade, drawn from a lognormal \\
Urgent share & 15\% & priority 1 of 3 \\
Promise to the complainant & 4 h / 1 day / 3 days & by priority \\
Dorm access & 40\% any time, 30\% morning, 30\% afternoon & somebody has to open the door \\
Cost of a walking minute & 0.05 & relative to one weighted hour late \\
\bottomrule
\end{tabular}

\caption{The assumptions behind the generated complaints. All of them live in
\texttt{src/instance.py} and all of them can be changed in one place.}
\label{tab:assume}
\end{table}

A few of these deserve a word. We measure time in working minutes, so the day
runs from $t=0$ at nine in the morning to $t=\NWorkingMinutes$ at half past
five, with the lunch hour simply left out of the clock. Complaints arrive at
each building at a steady average rate, which is why the number raised in a
given dorm on a given day follows a Poisson distribution. Priority sets two
things at once: how heavily lateness counts, and what we promised the
complainant. An urgent complaint is promised four hours, a normal one a full
working day, a minor one three days.

Access windows are the assumption that does the most work. We assume that
forty per cent of dorm complaints can be dealt with at any time, because
somebody is around or it is in a corridor, and that the rest are restricted to
either the morning or the afternoon. Offices and public buildings have no
restriction. As Section~\ref{sec:results} shows, this one assumption turns out
to explain almost all of the urgent complaints that still miss their deadline
even under the best plan we can find.

\section{What the model is trying to do}

The dispatcher chooses, for each technician $k$ and each pair of jobs $i$ and
$j$, whether $k$ travels directly from $i$ to $j$. Write $x_{ijk}=1$ when it
does and $0$ when it does not, and let $s_i$ be the time at which job $i$ is
started. A technician can only be given jobs in their own trade, so $x_{ijk}$
is fixed at zero whenever the trade of $j$ is not the trade of $k$. That one
restriction is what ties the assignment question to the routing question:
there is no point sending the nearest technician if the nearest technician is
a carpenter and the job is a leaking tap.

What we minimise is
\[
\sum_{i} w_i T_i \;+\; \alpha \sum_{i,j,k} c_{ij}\, x_{ijk},
\qquad
T_i = \max\{0,\; f_i - d_i\},
\]
where $f_i$ is when job $i$ was finished, $d_i$ is the deadline we promised,
$w_i$ is the weight of its priority class (ten for urgent, four for normal,
one for minor), and $c_{ij}$ is the walking time from $i$ to $j$. Lateness is
counted in hours and only in one direction: finishing early earns nothing.
The second term prices walking at $\alpha=\NAlpha$ per minute, which is small
on purpose. Walking is the tiebreaker, not the goal. If two plans close the
same complaints at the same times, we prefer the one with less walking; we
never accept a later repair in order to save a walk.

The constraints are the things that make a plan real rather than merely
cheap. Every job is done at most once and by exactly one technician. A
technician's route starts and ends at the Estate Office. Nobody can start a
job before its access window opens or after it shuts. Nobody can still be
working after the shift ends, and the walk back to the depot has to fit inside
the day too. Any job that does not fit is not lost: it stays in the queue and
appears again tomorrow, which is how a small daily shortfall turns into a
backlog.

That last point is worth being explicit about, because it changes the
arithmetic. Leaving a job undone looks free if you only look at today. It is
not free, so when the model weighs up whether to defer a job it charges that
job the lateness it would collect if it were done at the end of tomorrow.
Without that, the cheapest plan would always be to do nothing.

\section{Four ways to plan the day}

We built four methods, deliberately in increasing order of sophistication, so
that each one shows what the previous one was missing.

\subsection{First come, first served}

The baseline, and roughly what happens now. Take the complaints in the order
they were raised and give each one to the technician of the right trade who
has the least work so far. The route is simply the order they were assigned
in. No attention is paid to where anything is or to when it is due.

\subsection{Clarke and Wright's savings heuristic}

Start from the worst sensible plan: every job is its own trip, out from the
depot and straight back. That is obviously wasteful, so ask what you save by
doing jobs $i$ and $j$ on one trip instead of two:
\[
s_{ij} \;=\; c_{0i} + c_{0j} - c_{ij}.
\]
Compute this for every pair, sort the list from largest saving to smallest,
and work down it merging routes. A merge is only allowed if the two jobs are
currently on different routes, each is still at the end of its route, and the
merged route is still feasible: it fits in the shift and it respects every
access window.

This heuristic suits our campus unusually well, and it is worth seeing why.
Because the depot is on the old campus and nearly every complaint is on the
new campus, $c_{0i}$ is large for almost every job. That makes $s_{ij}$ large
for any pair of jobs on the same side of the road, so the savings list is
dominated by merges that bundle new-campus work into a single trip through the
tunnel. The heuristic discovers the tunnel for itself, without being told
about it.

The weakness of the method is equally clear: the savings formula contains
nothing about deadlines. It is minimising walking, and it will happily build a
beautifully efficient route that does the urgent job last.

\subsection{Priority insertion}

The third method attacks the problem from the other end. Start with empty
routes and insert jobs one at a time. For each unplaced job, find the cheapest
place to slot it into any technician's route, costing $\Delta_i$ extra minutes,
and then choose the job with the best ratio of what we avoid by doing it today
to what it costs to fit in. This is the bang-to-buck rule from the
profit-maximising tour heuristic, with urgency in place of revenue. When there
is more work than time, it fills the day with the jobs that matter and leaves
the minor ones behind, which is the behaviour we want.

\subsection{Improvement: 2-opt and Or-opt}

The fourth method takes the savings plan and tidies it. Three moves are tried
repeatedly until none of them helps.

\emph{2-opt} reverses a stretch of a route. On a map this removes crossings:
if a route doubles back on itself, flipping the segment between the two
crossing legs untangles it. \emph{Or-opt} lifts a run of one to three
consecutive jobs out of a route and reinserts it somewhere else. It is gentler
than 2-opt, which matters when access windows are tight, because it leaves the
order of everything else alone. The third move is Or-opt \emph{across} routes:
take a couple of jobs off one technician and give them to a colleague in the
same trade. This is the move that earns its keep, because savings merging
balances routes by geography and not by workload, so it regularly leaves one
technician with a packed day and another with a free afternoon.

Figures~\ref{fig:fifo} and~\ref{fig:imp} show what the difference looks like
for the three electricians on one ordinary day of \NDayTickets{} complaints.
The baseline crosses the tunnel repeatedly and wanders between campuses; the
optimised plan crosses once per technician and works each cluster out. Across
the whole crew that day, walking falls from \NDayFifoTravel{} minutes to
\NDayImpTravel{} minutes.

\figbox{routes_fifo}{First come, first served. Each numbered circle is a stop
in visiting order. The routes criss-cross and go back and forth through the
tunnel, because the order complaints arrived in has nothing to do with where
they are.}

\figbox{routes_improved}{The same three electricians, the same complaints,
after savings merging and improvement. Each technician now crosses the tunnel
once, works a cluster, and comes back.}

\section{How good are these plans really?}

A heuristic gives an answer with no guarantee attached, so the obvious
question is how far from the best possible answer it lands. For small
instances we can find out exactly, by branch and bound.

The search works on one technician's day. A node of the tree is a partial
route, the jobs done so far in order, and its children are the ways of adding
one more job. Searching every leaf is hopeless, so at each node we compute a
lower bound on the best that any completion of that partial route could
achieve, and if the bound is already worse than the best full route found so
far, the entire branch is discarded without being explored. It has been
fathomed.

Our bound has two parts, both chosen because they are easy to justify and
cheap to compute. Every job still to be done has to be reached from somewhere,
so the cheapest walk into it is a floor on what that leg can cost, and those
floors can be added up. And a job still to be done cannot possibly finish
before the current time plus its own repair time, so if that is already past
its deadline, that much lateness is unavoidable and belongs in the bound. Both
parts ignore the ordering, which is precisely why they never overstate the
true cost.

\begin{table}[h]\centering\small
\begin{tabular}{lrrrrrr}
\toprule
\textbf{Jobs} & \textbf{Optimum} & \textbf{Savings alone} & \textbf{Savings + improve} & \textbf{Gap} & \textbf{Nodes, warm} & \textbf{Nodes, cold} \\
\midrule
6 & 1.168 & 1.168 & 1.168 & 0.0\% & 321 & 340 \\
7 & 1.180 & 1.180 & 1.180 & 0.0\% & 1,516 & 1,580 \\
8 & 1.223 & 19.705 & 1.223 & 0.0\% & 4,051 & 9,702 \\
9 & 1.243 & 1.243 & 1.243 & 0.0\% & 8,959 & 22,078 \\
10 & 1.243 & 4.964 & 1.243 & 0.0\% & 48,069 & 125,603 \\
\bottomrule
\end{tabular}

\caption{The heuristics against the exact optimum, for one technician.
The last two columns count the nodes the same search had to explore when it
started from the heuristic's answer, and when it started from nothing and
explored jobs in arbitrary order.}
\label{tab:bb}
\end{table}

Table~\ref{tab:bb} gives the answer, and it is a good one: savings plus
improvement found the exact optimum on every instance we could solve, up to
\NBbMaxTickets{} jobs, with a worst gap of \NBbGapWorst\%. Savings on its own
does not, and the pattern of where it fails is instructive. At nine jobs it happens
to land exactly on the optimum; at \NBbSavWorstJobs{} jobs it is
\NBbSavWorstTimes{} times worse than it. That volatility is the price of a
method that never looks at deadlines: whether it does well depends entirely on
whether the geographically sensible route happens to put the urgent job early,
and that is luck rather than design.

The node counts in the last two columns make a second point, which is the one
the course makes about why heuristics matter inside exact methods. Starting
the search from the heuristic's answer cuts the work by \NBbNodeCutPct\%,
from \NBbNodesCold{} nodes to \NBbNodesWarm{}. A good incumbent lets the bound
start fathoming immediately instead of waiting for the search to stumble onto
a decent route by itself.

The honest limit of this section is that \NBbMaxTickets{} jobs for one
technician is nowhere near sixty jobs for nine technicians. What it tells us
is that the improvement moves are doing their job on problems we can verify,
which is the strongest evidence available without a commercial solver.

\section{Results}
\label{sec:results}

Table~\ref{tab:methods} compares the four methods over \NDays{} simulated
days. Every method sees exactly the same stream of complaints, generated from
the same random seed, so the differences are caused by the dispatching and
nothing else. Unfinished work rolls over, so a method that falls behind keeps
paying for it.

\begin{table}[h]\centering\small
\begin{tabular}{lrrrrrr}
\toprule
\textbf{Method} & \textbf{Objective} & \textbf{Late hours} & \textbf{Walk/tech/day} & \textbf{On time} & \textbf{Urgent on time} & \textbf{Avg waiting} \\
\midrule
First come first served & 3,603 & 3,241 & 26.8 min & 85.1\% & 37.0\% & 6.1 \\
Clarke-Wright savings & 2,766 & 2,564 & 15.0 min & 90.5\% & 55.5\% & 3.3 \\
Priority insertion & 358 & 106 & 18.6 min & 97.9\% & 85.7\% & 0.2 \\
Savings + 2-opt/Or-opt & 320 & 108 & 15.7 min & 98.0\% & 86.9\% & 0.6 \\
\bottomrule
\end{tabular}

\caption{The four methods over \NDays{} days, on an identical stream of
complaints. ``Late hours'' is the weighted lateness that the objective adds
up; ``avg waiting'' is the average number of complaints still unfinished at
the end of a day.}
\label{tab:methods}
\end{table}

Three things stand out.

Walking falls by about a third, from \NFifoTravel{} minutes per technician per
day to \NImpTravel. Across a crew of nine that is roughly
\NWalkSavedPerDay{} minutes of walking saved every day, or something close to
one and a half technician-days a week recovered and spent on repairs instead
of on corridors.

Service improves far more than walking does. The share of complaints closed
within the promised time rises from \NFifoSla\% to \NImpSla\%, and for urgent
complaints it roughly doubles, from \NFifoUrgent\% to \NImpUrgent\%. The
average wait for an urgent complaint falls from \NFifoUrgentWait{} hours to
\NImpUrgentWait{} hours, and the ninetieth percentile, which is what an
unlucky student actually experiences, falls from \NFifoUrgentNinety{} hours to
\NImpUrgentNinety.

The gain comes mostly from paying attention to urgency, not from saving
walking. Compare plain savings with priority insertion. Savings walks the
least of the three real methods, and still closes only \NSavUrgent\% of
urgent complaints on time. Priority insertion walks a little more and reaches
\NPriUrgent\%. Efficient routing is worth having, but it is the second most
important thing here.

\figbox{backlog}{Complaints still unfinished at the end of each day. The
baseline never clears its queue and spends the month carrying work forward.}

\figbox{sla}{Share of complaints closed within the promised time, by priority.
The difference between methods is almost entirely about urgent work: minor
complaints have three days to be dealt with and get done either way.}

\subsection{Why urgent complaints still miss, and what that means}

Even the best plan we found closes only \NImpUrgent\% of urgent complaints on
time, and it is worth asking what the remaining \NImpUrgentLate{} misses are.

They are almost entirely not the dispatcher's fault. Of the \NUrgentTotal{}
urgent complaints raised over the month, \NUrgentImpossible{} could never have
met the four hour promise no matter who was sent or in what order, because
their access window does not open until the afternoon and the repair itself is
long enough to run past the deadline. Even a technician already standing at
the door could not have made it. Of the \NImpUrgentLate{} urgent complaints
that the optimised plan closed late, \NImpUrgentUnavoidable{} are of exactly
that kind. That leaves \NImpUrgentAvoidable{} genuinely avoidable late urgent
complaint in thirty days.

The baseline, by comparison, misses \NFifoUrgentLate, of which
\NFifoUrgentAvoidable{} were avoidable.

This is the most useful result in the study, and it points somewhere other
than the algorithm. We are within one ticket of the best any dispatcher could
possibly do. The remaining failures are caused by the promise, not by the
plan, and no amount of further optimisation will touch them. Fixing them is a
policy decision: either the four hour promise should not apply to complaints
that can only be accessed in the afternoon, or urgent jobs need a way in
without the occupant, such as a dorm staff member with a key.

\section{What happens when it gets busy}

Thirty ordinary days is a soft test. The start of a term and the first week of
the monsoon both produce far more complaints than usual, so we repeated the
comparison with the arrival rate pushed from \NTicketsPerDay{} to
\num{85} complaints a day, which is about as much repair work as nine people
can physically absorb.

\begin{table}[h]\centering\small
\begin{tabular}{llrrrr}
\toprule
\textbf{Complaints} & \textbf{Method} & \textbf{On time} & \textbf{Urgent on time} & \textbf{Avg waiting} & \textbf{Unfinished at the end} \\
\midrule
58 a day & First come first served & 85.1\% & 37.0\% & 6.1 & 8 \\
58 a day & Savings + 2-opt/Or-opt & 98.0\% & 86.9\% & 0.6 & 0 \\
85 a day & First come first served & 53.9\% & 18.6\% & 113.5 & 208 \\
85 a day & Savings + 2-opt/Or-opt & 88.1\% & 83.9\% & 23.8 & 24 \\
\bottomrule
\end{tabular}

\caption{The same crew under normal load and under a surge.}
\label{tab:surge}
\end{table}

Under the surge the baseline collapses. Its backlog grows to an average of
\NSurgeHighFifoBacklog{} complaints, it ends the month with work still
outstanding, and urgent compliance falls to \NSurgeHighFifoUrgent\%. The
optimised dispatcher degrades, which is honest and expected, but it degrades
gracefully: \NSurgeHighImpSla\% of all complaints and \NSurgeHighImpUrgent\%
of urgent ones are still closed on time, and the backlog stabilises instead of
running away.

The reason is the deferral rule. When there is more work than time, somebody
has to wait, and the only question is who. The baseline defers whoever
happens to be at the back of the queue. The optimised dispatcher defers minor
complaints and protects urgent ones, which is why its urgent compliance barely
moves between the two scenarios while the baseline's halves.

\section{How many technicians, and of which trade?}

Because the model takes crew composition as an input, it can answer a staffing
question the Estate Office would actually care about.

\begin{table}[h]\centering\small
\begin{tabular}{lrrrrrr}
\toprule
\textbf{Crew} & \textbf{People} & \textbf{Objective} & \textbf{On time} & \textbf{Urgent on time} & \textbf{Walk/tech/day} & \textbf{Avg waiting} \\
\midrule
2 of each trade (6) & 6 & 2,718 & 90.3\% & 79.6\% & 22.8 min & 11.5 \\
3 of each trade (9) & 9 & 320 & 98.0\% & 86.9\% & 15.7 min & 0.6 \\
4 of each trade (12) & 12 & 293 & 98.1\% & 86.9\% & 11.8 min & 0.1 \\
4 electrical / 4 plumbing / 1 carpentry (9) & 9 & 2,356 & 93.4\% & 85.3\% & 15.6 min & 6.8 \\
\bottomrule
\end{tabular}

\caption{Different crews, all running the best method on the same complaints.}
\label{tab:crew}
\end{table}

Six technicians is not enough: the backlog runs away and urgent compliance
falls to \NCrewSixUrgent\%. Nine is enough. Going from nine to twelve is
almost pure waste, moving overall compliance from \NCrewNineSla\% to
\NCrewTwelveSla\% and urgent compliance not at all. Our recommendation would be
not to hire.

The fourth row was a surprise, and we have left it in because being wrong in
public is more useful than quietly deleting the experiment. Looking at the
workload, carpentry is clearly the quietest trade: it is twenty per cent of
complaints and the carpenters spend much of the day idle. So we tried keeping
the headcount at nine but shifting people to where the work is, giving four
electricians, four plumbers and one carpenter. It is dramatically worse:
the objective rises from \NCrewNineObj{} to \NCrewRebalObj. With only one
carpenter, carpentry becomes a bottleneck, and a single carpenter who falls
behind never catches up, because the backlog he carries makes the next day
worse. The idle time in the three-carpenter crew is not waste. It is the
buffer that absorbs a bad day, and removing it costs far more than it saves.

\section{Do these conclusions hold up?}

Two checks. The first is whether the result depends on the particular random
stream of complaints we happened to generate. We repeated the whole comparison
on \NSeedCount{} independent streams. Savings plus improvement scored between
\NSeedImpMin{} and \NSeedImpMax; the baseline scored between \NSeedFifoMin{}
and \NSeedFifoMax. The ranking of the methods is the same on every stream, and
the gap between best and baseline is far larger than the variation between
streams.

The second check is whether the conclusions depend on the assumptions we were
least sure about. Table~\ref{tab:sens} varies the two most arbitrary ones: the
price we put on a minute of walking, and the detour factor used to turn
straight-line distances into walking distances.

\begin{table}[h]\centering\small
\begin{tabular}{lllrrr}
\toprule
\textbf{Assumption} & \textbf{Value} & \textbf{Method} & \textbf{Objective} & \textbf{Walk/tech/day} & \textbf{On time} \\
\midrule
cost of a walking minute & 0.0 & First come first served & 3,603 & 26.8 min & 85.1\% \\
cost of a walking minute & 0.0 & Savings + 2-opt/Or-opt & 387 & 18.3 min & 97.8\% \\
cost of a walking minute & 0.05 & First come first served & 3,603 & 26.8 min & 85.1\% \\
cost of a walking minute & 0.05 & Savings + 2-opt/Or-opt & 320 & 15.7 min & 98.0\% \\
cost of a walking minute & 0.25 & First come first served & 3,603 & 26.8 min & 85.1\% \\
cost of a walking minute & 0.25 & Savings + 2-opt/Or-opt & 319 & 15.3 min & 97.9\% \\
walking detour factor & 1.2 & First come first served & 3,566 & 25.1 min & 85.1\% \\
walking detour factor & 1.2 & Savings + 2-opt/Or-opt & 319 & 14.9 min & 98.0\% \\
walking detour factor & 1.3 & First come first served & 3,603 & 26.8 min & 85.1\% \\
walking detour factor & 1.3 & Savings + 2-opt/Or-opt & 320 & 15.7 min & 98.0\% \\
walking detour factor & 1.4 & First come first served & 3,684 & 28.5 min & 85.0\% \\
walking detour factor & 1.4 & Savings + 2-opt/Or-opt & 334 & 16.7 min & 98.0\% \\
\bottomrule
\end{tabular}

\caption{Varying the two most arbitrary assumptions in the model.}
\label{tab:sens}
\end{table}

Neither changes the conclusion, though the first is more interesting than we
expected. Setting $\alpha$ to zero tells the model not to care about walking
at all, and walking rises only from \NImpTravel{} to \NAlphaZeroTravel{}
minutes per technician per day. Most of the saving survives even when nobody
is asking for it, because walking less is largely a by-product of planning the
day sensibly rather than something the objective has to insist on. The rest of
it, about \NAlphaWorthMinutes{} minutes per technician per day, is what
explicitly pricing walking buys. Pushing the detour factor from \num{1.2} to
\num{1.4} scales everyone's walking up together and leaves the comparison
between methods untouched.

\section{What we would tell the Estate Office}

\begin{itemize}
\item Plan the day in the morning rather than working the queue as it
  arrives. On our figures this is worth about a third of the crew's walking
  and roughly doubles the share of urgent complaints closed on time, with no
  extra staff and no new equipment.
\item Bundle by campus. The single biggest avoidable cost is crossing the
  tunnel more often than necessary. A technician should cross, work the
  cluster, and come back.
\item Do not hire more technicians. Nine is enough for the load we modelled,
  and the twelfth person buys almost nothing. If complaint volumes rise to
  surge levels permanently, revisit this.
\item Do not rebalance the trades to match average workload. The quiet trade
  is quiet on average and not on a bad day, and cutting it to one person makes
  the whole system fragile.
\item Change the promise, or change the access. Nearly all the urgent
  complaints that are still closed late are ones where the student is only
  available in the afternoon and the repair is too long to finish by the
  deadline. That cannot be fixed by better scheduling. It needs either a
  different promise for access-restricted complaints or a way of getting into
  the room without the occupant.
\end{itemize}

\section{What we have not done}

The model assumes repair times are known when the day is planned. In reality
a job that was expected to take twenty minutes sometimes takes ninety, and a
plan built on the wrong estimate will slip. A sensible extension would be to
re-plan at midday with whatever has actually happened by then.

We plan each day in the morning and do not react to anything that happens
during it. A genuinely urgent complaint raised at eleven o'clock waits for
tomorrow's plan in our model, which is not how an Estate Office would
behave. Allowing jobs to be inserted into a running route is the obvious next
step, and the cheapest-insertion machinery is already there to do it.

Our complaint data is invented. The assumptions are reasonable and they are
all in one file, but they are still assumptions. Of all of them, the access
window mix is the one that most deserves to be checked against reality,
because Section~\ref{sec:results} shows it is what limits the result.

Finally, we verify against an exact optimum only for a single technician with
up to \NBbMaxTickets{} jobs. Verifying the full nine-technician problem would
need a proper mixed-integer solver.

\section{The code}

Everything is in the repository that accompanies this report. It is plain
Python with no external libraries, so it runs anywhere with
\texttt{python3} and nothing installed. \texttt{python3 run.py} rebuilds the
walking network, runs every experiment, redraws the figures, writes the Excel
workbook and regenerates the numbers used in this document.

\begin{itemize}\small
\item \texttt{src/campus.py} builds the walking network from the map data.
\item \texttt{src/instance.py} generates complaints and holds every assumption.
\item \texttt{src/dispatch.py} contains the four planning methods.
\item \texttt{src/exact.py} is the branch and bound solver.
\item \texttt{src/experiment.py} runs the multi-day simulation.
\item \texttt{src/study.py} runs everything and writes the outputs.
\end{itemize}

\end{document}
